\subsection{Условия постоянства и монотонности функций.}

Все доказательства в этом параграфе опираются на формулу Лагранжа:
\[ f(x_2) - f(x_1) = f'(\xi)(x_2 - x_1), \quad \text{где } \xi \in (x_1, x_2). \]

\subsubsection*{Критерий постоянства функции}

\begin{tcolorbox}[title=Теорема (Критерий постоянства)]
    Пусть функция $f(x)$ дифференцируема на интервале $(a, b)$.
    Функция является постоянной на $(a, b)$ тогда и только тогда, когда её производная тождественно равна нулю:
    \[ f(x) = \text{const} \iff \forall x \in (a, b): f'(x) = 0 \]
\end{tcolorbox}

\textbf{Доказательство:}
\begin{enumerate}
    \item ($\Rightarrow$) Необходимость.
    Если $f(x) = C$, то производная константы равна нулю: $(C)' = 0$.
    
    \item ($\Leftarrow$) Достаточность.
    Пусть $f'(x) = 0$ во всех точках интервала.
    Возьмем любые две точки $x_1, x_2 \in (a, b)$ (пусть $x_1 < x_2$).
    По теореме Лагранжа на отрезке $[x_1, x_2]$:
    \[ f(x_2) - f(x_1) = f'(\xi)(x_2 - x_1). \]
    Так как по условию $f'(\xi) = 0$, то:
    \[ f(x_2) - f(x_1) = 0 \implies f(x_2) = f(x_1). \]
    Значения функции в любых точках совпадают, значит $f(x) = \text{const}$.
\end{enumerate}

\subsubsection*{Условия монотонности}

\begin{tcolorbox}[title=Теорема (Критерий монотонности)]
    Пусть $f(x)$ дифференцируема на $(a, b)$.
    \begin{enumerate}
        \item $f(x)$ неубывает ($\nearrow$) на $(a, b) \iff \forall x \in (a, b): f'(x) \ge 0$.
        \item $f(x)$ невозрастает ($\searrow$) на $(a, b) \iff \forall x \in (a, b): f'(x) \le 0$.
    \end{enumerate}
\end{tcolorbox}

\textbf{Доказательство (для неубывания):}
\begin{enumerate}
    \item ($\Rightarrow$) Необходимость.
    Пусть $f(x)$ неубывает. Тогда для $x > x_0$ имеем $f(x) \ge f(x_0)$.
    Разностное отношение $\frac{f(x) - f(x_0)}{x - x_0} \ge 0$.
    Переходя к пределу при $x \to x_0$, получаем $f'(x_0) \ge 0$.
    
    \item ($\Leftarrow$) Достаточность.
    Пусть $f'(x) \ge 0$. Возьмем произвольные $x_1 < x_2$.
    По теореме Лагранжа:
    \[ f(x_2) - f(x_1) = \underbrace{f'(\xi)}_{\ge 0} \cdot \underbrace{(x_2 - x_1)}_{> 0} \ge 0. \]
    Следовательно, $f(x_2) \ge f(x_1)$, то есть функция не убывает.
\end{enumerate}

\subsubsection*{Условия строгой монотонности}

\begin{tcolorbox}[title=Достаточное условие строгого возрастания]
    Если $f'(x) > 0$ на интервале $(a, b)$, то $f(x)$ \textbf{строго} возрастает на этом интервале.
\end{tcolorbox}

\textbf{Доказательство:}
Аналогично предыдущему: $x_2 > x_1 \implies x_2 - x_1 > 0$.
Так как $f'(\xi) > 0$, то их произведение строго больше нуля:
\[ f(x_2) - f(x_1) > 0 \implies f(x_2) > f(x_1). \]

\textbf{Важное замечание:}
Условие $f'(x) > 0$ является только \textit{достаточным}, но не необходимым.
Функция может строго возрастать, даже если производная обращается в ноль в отдельных точках.
\textit{Пример:} $y = x^3$. Она строго возрастает на всей $\mathbb{R}$, но в точке $x=0$ её производная $y'=3x^2$ равна нулю.
